\documentclass{article}

\usepackage[utf8]{inputenc}
\usepackage[T1]{fontenc}
\usepackage{helvet}
\usepackage{geometry}
\usepackage{xcolor}
\usepackage{eso-pic}
\usepackage{fancyhdr}
\usepackage{parskip}
\usepackage{microtype}
\usepackage[english, ukrainian]{babel}
\usepackage{url}
\usepackage{amsmath}
\usepackage{amssymb}
\usepackage{amsthm}
\usepackage{longtable}
\usepackage{booktabs}
\usepackage{hyperref}
\usepackage{titlesec}

\definecolor{nato}{HTML}{293757}
\definecolor{footergray}{gray}{0.65}

\geometry{a4paper,left=3cm,right=3cm,top=3cm,bottom=3cm}
\renewcommand{\familydefault}{\sfdefault}
\setcounter{secnumdepth}{3}

\titleformat{\section}
  {\color{nato}\Huge\bfseries}{\thesection.}{1em}{}
\titlespacing*{\section}{0pt}{30pt}{12pt}

\titleformat{\subsection}
  {\color{nato}\LARGE\bfseries}{\thesubsection.}{1em}{}
\titlespacing*{\subsection}{0pt}{20pt}{8pt}

\titleformat{\subsubsection}
  {\color{nato}\Large\bfseries}{\thesubsubsection.}{1em}{}
\titlespacing*{\subsubsection}{0pt}{10pt}{6pt}

\fancyhf{}
\fancyhead[R]{\small\color{footergray} 31 березня 2026}
\fancyfoot[L]{\small\color{footergray} Технічні вимоги ERP/1. Політики}

\fancyfoot[R]{\small\color{footergray}\thepage}

\newcommand\HUGE[1]{\fontsize{40}{40}\selectfont #1}

\renewcommand{\headrulewidth}{0pt}
\renewcommand{\footrulewidth}{0.4pt}
\renewcommand{\footrule}{\color{footergray}\hrule width \textwidth height 0.4pt}

\begin{document}

\pagestyle{fancy}

\begin{titlepage}
  \AddToShipoutPictureBG*{\AtPageLowerLeft{\color{nato}\rule{\paperwidth}{\paperheight}}}
  \color{white}\thispagestyle{empty}\vspace*{4cm}\centering
  {\Huge  \textbf{ERP/1: Авторизація} \par} \vspace{2cm}
  {\HUGE  \textbf{Політики} \par} \vspace{2.5cm}
  {\LARGE \textbf{Технічних вимог, Техноробочого проєкту \\ і Технічної документації} \par} \vspace{2.5cm}
  \vfill
  {\large 2026 \copyright\ Системи Електронного Урядування \par}
\end{titlepage}

\newpage

\begin{center}
  {\Huge\color{nato}\bfseries Зміст\par}
\end{center}
\vspace{40pt}
\tableofcontents

\begin{abstract}
  У статті подаються універсальні засади політик проектів для
  державних ІТ підприємств, типовий зміст технічних вимог з повною таксономією підрозділів та
  структурою технічного завдання відповідно до ДСТУ.
  Показані принципи, як політики визначають вимоги, що узгоджують технічне завдання та посилання на ключові
  державні стандарти документування ДСТУ. Матеріал призначений для архітекторів, аналітиків та розробників,
  які працюють над державними та комерційними проєктами в Україні.
\end{abstract}

\newpage

\section{Політики}

Політики є первинним і незмінним документом будь-якого проєкту.
Вони визначають правила гри до того, як формулюються вимоги чи обираються технології.

\subsection*{Умовні скорочення та визначення}

Терміни та скорочення, що використовуються в документі:

\begin{longtable}{p{4cm}p{10cm}}
\toprule
\textbf{Термін} & \textbf{Значення} \\
\midrule
ТВ & Технічні вимоги \\
ТЗ & Технічне завдання \\
ТД & Технічна документація \\
\bottomrule
\end{longtable}

\subsection*{Загальні принципи}

\begin{enumerate}
\item Політики є первинними документами.
\item Усі рішення в проєкті приймаються виключно на основі затверджених політик.
\item Технічні вимоги формуються виключно як похідні від політик.
\item Технічне завдання є єдиним документом, де дозволяється вказувати конкретні продукти.
\item Архітектура повинна залишатися агностичною до моменту затвердження ТЗ.
\item Загальні нефункціональні вимоги повинні міститися в цих політиках.
\item Технічні вимоги не повинні містити реквізитних частин сутностей і документів.
\item Технічні вимоги повинні містити перелік сутностей.
\item Технічні вимоги повинні містити перелік процесів.
\item Технічні вимоги повинні містити перелік прав та превілеїв.
\item Технічні вимоги доменної області повинні містити перелік процесів і їх кроків.
\item Технічні завдання повинні містити BPMN діаграми, верифіковані аналітиками.
\item Технічні завдання повинні містити реквізитні частини сутностей і документів.
\item Технічне завдання повинно містити перелік ролей і ABAC правил.
\item Технічне завдання повинно містити перелік словників.
\item Технічне завдання повинно містити перелік протоколів взаємодії.
\item Технічне завдання повинно містити перелік шаблонів.
\item Будь-яка інформація, що дублюється в похідних документах повинна міститись тут.
\item Технічна документація повинна визначатись стандартами ISO і ДСТУ.
\item Технічне завдання повинно містити продукти, які підтримують стандарти ISO або ДСТУ.
\end{enumerate}

\newpage
\section{Технічні вимоги}

\subsection{Перелік технічних вимог наведених в політиках}

\begin{itemize}
\item Функціональні вимоги до словникової системи
\item Функціональні вимоги до набору облікових сутностей і їх операцій
\item Функціональні вимоги до набору валідацій
\item Функціональні вимоги до рольової моделі
\item Функціональні вимоги до інтерфейсу користувача
\item Функціональні вимоги до системи процесів
\item Нефункціональні вимоги до архітектури
\item Нефункціональні вимоги до безпеки
\item Нефункціональні вимоги до потужності та ємності
\item Нефункціональні вимоги до системи логуання
\item Нефункціональні вимоги до інтерфейсу користувача
\item Нефункціональні вимоги до контролю якості
\item Нефункціональні юридичні вимоги
\end{itemize}

\newpage
\subsection{Повний зміст технічних вимог}
Структура технічних вимог узагальнена для будь-якого високотехнологічного державного ПЗ.

\begin{itemize}
\item Умовні скорочення та визначення
\item Загальні відомості
  \begin{itemize}
  \item Передумови
  \item Питання, що вирішуються
  \item Вимоги законодавства та міжнародних стандартів (ДСТУ, ISO, NIST, ITU-T)
  \end{itemize}
\item Призначення та цілі впровадження
\item Класифікація вимог
  \begin{itemize}
  \item 3.1. Функціональні вимоги
    \begin{itemize}
    \item 3.1.1. Опис довідників
    \item 3.1.2. Опис реєстрових сутностей та їх станів
    \item 3.1.3. Опис бізнес-процесів (BPMN)
    \item 3.1.4. Вимоги до протоколів взаємодії
    \item 3.1.5. Опис прав та привілеїв (ABAC/RBAC/ACL)
    \item 3.1.6. Опис валідацій
    \item 3.1.7. Шаблони сповіщень
    \item 3.1.8. Вимоги до серіалізації
    \item 3.1.9. Вимоги до транспорту
    \end{itemize}
  \item 3.2. Нефункціональні вимоги
    \begin{itemize}
    \item 3.1.1. Вимоги до архітектури
    \item 3.2.1. Вимоги до безпеки
    \item 3.2.2. Вимоги до потужності і ємності
    \item 3.2.4. Вимоги до функціональності логування
    \item 3.2.5. Адміністративні вимоги
      \begin{itemize}
      \item 3.2.5.1. Настанова до видів тестування
      \item 3.2.5.2. Вимоги до настанов адміністратора
      \item 3.2.5.3. Вимоги до документування
      \item 3.2.5.4. Технологічний стек (агностичний опис)
      \end{itemize}
    \item 3.2.6. Загальні вимоги до документації та артефактів
    \item 3.2.7. Юридичні вимоги
    \end{itemize}
  \end{itemize}
\end{itemize}

\newpage
\section{Технічне завдання}
Структура технічного завдання відповідає \textbf{ДСТУ 3973-2000} «Система розроблення та постачання
продукції на виробництво. Правила виконання науково-дослідних робіт. Загальні
положення» та \textbf{ДСТУ 3974-2000} (для конструкторських робіт).

\subsection{Зміст технічного завдання}
ДСТУ 3973-2000 визначає зміст технічного завдання.

\begin{enumerate}
\item Мотивація, мета, виконавці
\item Класифікація вимог (функціональні, нефункціональні, технічні, адміністративні, ергономічні,
      політичні, юридично-правові, документальні, архітектурні)
\item Принципові вимоги (відповідність ДСТУ, NIST, ISO, ITU-T, ДССЗІ)
\item Функціональні вимоги
\item Ергономічні вимоги
\item Технічні вимоги
\item Адміністративні вимоги
\item Політичні вимоги
\item Юридично-правові вимоги
\item Документування (перелік документів)
\item Етапи виконання
\item Результати
\item Порядок передачі
\end{enumerate}

\subsection{Вимоги до Технічних завдань}
\subsubsection{Правила-рекомендації до ТЗ}

\begin{itemize}
    \item \textbf{Продовження технічних вимог.} Технічне завдання має бути продовженням технічних вимог. Кожна вимога, або група вимог визначена на рівні технічних вимог, повинна бути відображена в ТЗ, деталізована та пов’язана з конкретною функціональністю. Має забезпечуватися простежуваність від початкової потреби до реалізації, тестування та приймання результату + зв’язок із сутностями.
    \item \textbf{Людиноцентричний підхід.} Розробка має ґрунтуватися на людиноцентричному підході.
    \item \textbf{Підтверджені потреби користувачів.} Функціональність повинна створюватися для вирішення підтверджених потреб конкретних груп користувачів. Користувацькі сценарії, прототипи та дизайн-рішення мають перевірятися з реальними представниками відповідних ролей до завершення розробки.
    \item \textbf{Контекст бізнес-процесів.} Функціональність необхідно описувати в контексті бізнес-процесів. Для кожного рішення має бути зрозуміло, який процес або його конкретний крок воно підтримує, хто бере участь у процесі, з чого він починається та який результат повинен бути отриманий.
    \item \textbf{Розкриття через Use Case.} Ключова функціональність має розкриватися через use case. Сценарії використання повинні показувати взаємодію користувача, системи та інших учасників для досягнення визначеної мети. Рівень деталізації має залежати від складності, ризику та критичності функціональності.
    \item \textbf{Альтернативні та негативні сценарії.} Use case повинні охоплювати не лише основний успішний сценарій. Необхідно враховувати суттєві альтернативні, негативні та виняткові сценарії, зокрема помилки користувача, відсутність даних, недоступність інтеграцій, повторне виконання дій та відновлення процесу.
    \item \textbf{Визначення бізнес-правил у BPMN.} У межах кожного use case мають бути явно визначені бізнес-правила, які керують відповідним сценарієм. До них належать умови перевірки даних, маршрутизації процесу, застосування строків, розрахунків, ролей, обмежень і прийняття рішень. Такі правила не повинні самостійно визначатися розробником під час реалізації. За потреби для них зазначаються джерело, умови застосування та винятки.
    \item \textbf{Явне визначення бізнес-правил.} Бізнес-правила мають бути визначені явно. Правила перевірки, маршрутизації, строків, розрахунків, ролей і прийняття рішень не повинні формуватися розробником самостійно під час реалізації. Для критичних правил мають бути визначені їх джерело, умови застосування та винятки.
    \item \textbf{Критерії приймання.} До початку розробки потрібно визначити критерії приймання. Для кожної функціональності має бути зрозуміло, який результат очікується та за якими об’єктивними ознаками буде перевірятися його відповідність вимогам.
    \item \textbf{Definition of Done.} Окремо має бути визначений Definition of Done. Він повинен встановлювати спільний мінімальний стандарт завершеності: реалізація функціональності, виконання критеріїв приймання, тестування, документування, перевірка безпеки, доступності, інтеграцій та усунення критичних дефектів.
    \item \textbf{Перевірка відповідності та цінності.} Результат необхідно перевіряти як на відповідність ТЗ, так і на реальну цінність. Важливо підтвердити не лише те, що функціональність реалізована відповідно до опису, а й те, що вона працює в реальному процесі, вирішує початкову проблему користувача та може бути безпечно впроваджена й масштабована.
    \item \textbf{Наскрізна простежуваність (Traceability).} Має бути забезпечена наскрізна простежуваність від політик безпеки до технічних вимог, технічного завдання, структури бази даних, вихідного коду та сценаріїв тестування. Кожен елемент реалізації системи повинен мати чітку прив'язку до первинної вимоги для проходження міжнародного аудиту.
\end{itemize}

\newpage
\section{Посилання на стандарти ДСТУ}

\begin{longtable}{p{4cm}p{9cm}}
\toprule
\textbf{Стандарт} & \textbf{Назва та призначення} \\
\midrule
ДСТУ 3973-2000 & Система розроблення та постачання продукції. Правила виконання науково-дослідних робіт. Загальні положення (основа структури ТЗ) \\
ДСТУ 3974-2000 & Правила виконання науково-конструкторських робіт \\
ДСТУ 42010:2018 & Інженерія систем і програмних засобів. Опис архітектури \\
ДСТУ 62264:2019 & Інтеграція систем управління підприємством і виробництвом \\
ДСТУ 3008:2015 & Документація на розроблення ПЗ \\
ДСТУ 3396.0-96, 3396.1-96, 3396.2-97 & Комплекс стандартів на уніфіковану систему документації \\
ДСТУ 4541, ДСТУ 28147 & Криптографічні стандарти України \\
ДСТУ 19.403--19.781, ДСТУ 34.003, 34.201, 34.601--34.602 & Застарілі, але обов’язкові для сумісності стандарти документування \\
Постанова КМУ № 205 & Порядок використання засобів інформатизації від 21.02.2025 \\
\bottomrule
\end{longtable}

\newpage
\section{Функціональні вимоги}

\subsection{Вимоги до процесів системи}
Система повинна буди ізоморвна моделі FSM (Finite State Machines), BPMN (ISO 19510), або Мережам Петрі,
тобто повинна підтримувати парелальні і послідовні обчислення (погодження) необхідні для орагіназіції
документообігу. Визначення процесу, як сутності, що складається з наступних частин (як приклад такої ізоморфної системи):

\begin{itemize}
    \item Контекст процесу у вигляді переліку документів
    \item Перелік станів процесу, кожен з яких є робочим місцем користувача або сервісом
    \item Переходи між станами що є функціями з логікою
    \item Початковий і кінцевий стани процесу
    \item Історія кроків процесу з документами
\end{itemize}

Кожен користувач визначається як сутність, що містить наступні реквізити:
\begin{itemize}
    \item Сертифікат доступу КЕП X.509
    \item ЄДРПО підприємства
    \item Посада, ПІБ, табельний номер
    \item Визначені ролі
\end{itemize}

Система бізнес-процесів повинна формально визначати:
\begin{itemize}
    \item Перелік користувачів і структури підприємства
    \item Перелік ролей
    \item Перелік реєстрових сутностей, їх реквізитів і операцій над ними
    \item Перелік документів і їх реквізитів
    \item Перелік інформаційних повідомлень і їх реквізитів
    \item Перелік бізнес-процесів
\end{itemize}

\newpage
\subsection{Вимоги до викликів API}
Для початку розробки аналітик замовника передає вендору зазначені верхньорівневі вимоги до методу.
Після завершення розробки вендор надає опис методу за затвердженим шаблоном відповідно до настанови.
Документація по опису методів API повинна бути створена згідно з затвердженими шаблонами.

\begin{itemize}
    \item призначення методу: яку дію він має виконувати;
    \item формулювання методу: метод та ендпойнт;
    \item які параметри необхідні для виклику методу — початкові умови для застосування методу (параметри мають бути представлені у формі таблиці із зазначеним типом та прикладом формулювання);
    \item відповідь сервера у формі структурованих даних — інформація у форматі серіалізації (зразок відповіді описують у формі коду);
    \item обробка помилок: опис ситуації, коли виконання дії неможливе з певних причин (перелік помилок мають бути представлені у формі таблиці з розшифровкою значення);
    \item додаткова інформація: обмеження доступу, використання стандартів тощо (може бути представлена у формі довільного тексту в окремому розділі).
\end{itemize}

\subsection{Вимоги до інтерфейсу користувача}

\subsubsection{Обробка помилок}
Інтерфейс користувача має представляти чіткі та інформативні повідомлення про помилки із
зазначенням їхньої причини та шляху виправлення. Інтерфейс користувача має надавати можливість
повторної спроби операцій після усунення факторів, які спричинили помилку.

\subsubsection{Загальні вимоги}
Інтерфейс користувача має відповідати наступним загальним вимогам:

\begin{itemize}
    \item \textbf{Відповідність стандартам вебдоступності (WCAG):} Забезпечення доступності для
           користувачів з обмеженими можливостями (наприклад, підтримка клавіатурної навігації,
           альтернативний текст для зображень, достатній контраст кольорів). Відповідність принаймні рівню AA стандартів WCAG 2.1.
    \item \textbf{Наслідування дизайн-коду компанії:} Інтерфейс має відповідати існуючому дизайн-коду
           компанії (стиль, кольори, шрифти, логотипи тощо). Використання компонентів UI з бібліотеки
           компонентів компанії (якщо така є для певного функціоналу). Забезпечення консистентності
           візуального стилю між різними частинами системи.
\end{itemize}

\addtocontents{toc}{\protect\newpage}
\newpage
\section{Нефункціональні вимоги}

\subsection{Вимоги до архітектури}

Перший підрозділ описує загальну архітектуру ERP/1, та вимоги до розробників додаткової функціональності.
Ця глава визначає карту вимог та описує загальний перелік вимог, який випливає з логіки розгортання
архітектурних рівнів ISO-42010.

\subsubsection{Визначення термінів і призначення документа}

Архітектура тут і надалі буде означати формальний структурний опис компонент і регламентів взаємодії з розширеною деталізацією прикладного рівня.

Призначення документа:
\begin{itemize}
  \item ознайомити виконавця з поточною архітектурою ERP/1;
  \item визначити технології, які існують й використовуються в ERP/1;
  \item визначити вимоги та обмеження реалізації.
\end{itemize}

\subsubsection{Рівні архітектури компонентів системи}

За основу взята архітектурна методологія ISO/IEC/IEEE 42010 та фреймворк Закмана.

\begin{enumerate}
  \item Компоненти архітектури ERP/1 є цілком визначеними та узгодженими.
  \item Контрактори, які реалізують проєкт нової функціональності ERP/1 повинні зазначити на
        наступних архітектурних рівнях які модифікації потребує система:
    \begin{itemize}
      \item Сутності (Resources, ER-діаграми);
      \item Форми (Forms, містять додаткову інформацію для UI);
      \item Валідації (Validations);
      \item Процеси (Послідовність використання API, BPMN діаграми);
      \item Довідники (Dictionaries);
      \item Права (Scopes);
      \item Функції сервісів (публічні та приватні) (Functions, API, RPC);
      \item SQL-таблиці;
      \item KV-простори;
      \item Cache-простори;
      \item Правила обмеження швидкості засобами KONG (Rate limiters).
    \end{itemize}
  \item Перелік модифікацій повинен надаватися в формі, де кожна зміна зазначена у відповідному
        документі-вимозі до оформлення кожного архітектурного рівня. В цьому контексті вимоги
        до опису архітектури є кореневим документом-вимогою який описує зміни до архітектури ERP/1
        на верхньому архітектурному рівні.
\end{enumerate}

\newpage
\subsubsection{Маніфест відповідності архітектурі ERP/1}

Мотиваційна частина: мета документу.

\begin{enumerate}
  \item Стандартизація та уніфікація функцій ПЗ повинна бути забезпечена за рахунок використання сучасних інструментальних програмних засобів, які підтримують єдину технологію проектування і розробки функціонального, інформаційного та програмного забезпечення.
  \item ПЗ в цілому та інші програмні компоненти повинні відповідати основним міжнародним та національним угодам і стандартам в галузі інформаційних технологій.
  \item Необхідно використовувати вже наявні у Замовника системи та програмне забезпечення.
  \item Версії систем і підсистем продуктів повинні відповідати версіям, що використовуються.
  \item Нове технічне забезпечення (наприклад, підсистеми) має відповідати наступним вимогам:
    \begin{itemize}
      \item безкоштовність нової системи/підсистеми, її залежностей, і їх ліцензій;
      \item наявність підтримки розробника (1 рік від останньої версії);
      \item програмне забезпечення з відкритим вихідним кодом (open source).
    \end{itemize}
  \item Вимоги стосуються безпосередньо і програмного забезпечення, розробку якого сплачує Замовник.
  \item Будь-яке інше програмне забезпечення повинно відповідати вимогам.
  \item Вимоги до горизонтального масштабування (за необхідності). Шардінг, контейнеризація, оркестрація.
  \item Балансування навантаження між декількома екземплярами сервісів (подами).
  \item Використання черг для часомістких операцій.
  \item Мінімізація розміру інформаційних повідомлень і оптимізація трафіку.
  \item Вимоги до телеметрії, кількість операцій за секунду, латенсі, помилки, статуси.
  \item Вимоги до потужності і ємності.
  \item Вимоги до логування.
  \item Можливість проведення автоматичного тестування з відключеними сервісами КЕП.
  \item Мінімізація вартості володіння.
\end{enumerate}

\newpage
\subsubsection{Структура продуктів ERP/1}

\paragraph{Корпоративні продукти}
\begin{itemize}
    \item \textbf{Освіта} --- Автоматизована інформаційна система управління освітнім процесом ISO 21001, навчальною діяльністю та адмініструванням закладів освіти з підтримкою очної, дистанційної та змішаної форм навчання.
    \item \textbf{Здоров'я} --- Медична інформаційна система HL7 для автоматизації обліку медичних послуг та управління медичною інформацією в електронному вигляді.
    \item \textbf{Документи} --- Автоматизація роботи з електронними документами у вигляді CRM системи згідно інструкції з діловодства №40 від 2024 року.
    \item \textbf{Облік} --- Комплексна ERP система автоматизації бухгалтерського обліку, кадрового діловодства та розрахунку заробітної плати для підприємств і державних установ.
    \item \textbf{Склад} --- Автоматизована система управління складським господарством та матеріально-технічним забезпеченням підприємств і державних установ.
    \item \textbf{Реєстри} --- Універсальна платформа для створення, ведення та управління державними, відомчими та корпоративними інформаційними реєстрами будь-якого масштабу.
\end{itemize}

\paragraph{Системні продукти (Платформа)}
\begin{itemize}
    \item \textbf{Сертифікати} --- Центр сертифікації реєстраційних посвідчень і їх ABAC реквізитної атрибутики для інфраструктури PKI.
    \item \textbf{Тунелі} --- Високопродуктивна система побудови захищених децентралізованих тунелів (L2/L3) на базі віртуальної машини Erlang/OTP.
    \item \textbf{Директорія} --- PKI-aware LDAP з підтримкою X.500.
    \item \textbf{Комунікатор} --- TLS X.690 DER сервер v1+, v2 і v3 протоколів і ПК КЗІ X.509 v1 протоколу BUDDHA.
    \item \textbf{Процеси} --- Процесний BPMN рушій ISO 19510 для бізнес-логіки з TLS X.690 DER інтерфейсом.
    \item \textbf{Брокер} --- MQTT v5.0 сервер персистентних топіків стандарту ISO/IEC 20922:2016, що використовується в продуктах SYNRC.
    \item \textbf{Сховище} --- Бібліотека KVS, що реалізує SNIA інтерфейс для NVMe дисків, сумісна з RocksDB та CEPH.
    \item \textbf{Фреймворк} --- N2O/NITRO фреймворк для розробки веб-додатків.
    \item \textbf{Форми} --- Бібліотека X-FORM для відображення реквізитної інформації та автоматичної генерації клієнтських і серверних валідацій.
    \item \textbf{Компілятор} --- ASN.1 X.680 компілятор для мов програмування ANSI C99 та Apple Swift.
    \item \textbf{Контроль} --- Бібліотека ABAC для контролю доступу на рівні реквізитної інформації документів та їх сертифікатів.
\end{itemize}

\newpage
\subsubsection{Таксономія системних компонент}

Перелік категорій компоненти згідно з класифікацією по їх функціональним категоріям (в дужках зазначені
приклади які згодом можливо зафіксуються в технічному завданні).

\begin{itemize}
  \item Data Storages (Приклади: MongoDB, SQL, S3, Cassandra, Riak, tikv);
  \item PubSub Buses (Приклади: Kafka, MQTT, AMQP);
  \item Application Services (Приклади: D3, ERP/1, CRM);
  \item Facade Management (Приклади: Traefik / KONG); OSI Layer 4—7;
  \item Virtual Machines (Приклади: JVM, CLR, BEAM);
  \item External Systems (Трембіта: ДССУ, Дія, ПФУ, Мінюст).
\end{itemize}

Види Data Storages:
\begin{itemize}
  \item KV. Сховище з єдиним простором ключів для зберігання подій і документів.
  \item SQL. Реляційне сховище для обліку бізнес-сутностей.
  \item S3. Сховище медіа-об'єктів (документів, файлів, тощо) з Amazon S3 інтерфейсом.
  \item FS. IPFS, NTFS, SAMBA, GlusterFS, ZFS.
\end{itemize}

Види Pub/Sub Buses:
\begin{itemize}
  \item MQTT. Брокер повідомлень, основа Service Oriented Architecture (SOA). ISO.
  \item AMQP. Теж. ISO. Не використовується зараз.
  \item XMPP. Теж. ISO.
  \item Kafka. Non-ISO.
  \item NSQ. Теж. Non-ISO.
\end{itemize}

Application Services:

\begin{itemize}
  \item Anti Fraud. Сервіс блокування зловмисних дій.
  \item One Time Password. Сервіс верифікації через одноразові паролі.
  \item Authentication and Authorization. Сервіс аутентифікації і авторизації.
  \item Integration Layer. Основні сервіси і реєстри.
  \item Data Layer. Безпосередній доступ до реплік PostgreSQL.
  \item Digital Signature. Сервіс елекронного підпису.
  \item Master Citizen Index (MCI). Сервіс ізоляції персональних даних користувачів.
  \item Extract, Transform, Load (ETL). Зовнішні конектори.
\end{itemize}

Web Frameworks:

\begin{itemize}
  \item Основані на шаблонізаторах (PHP, Django DTL, Elixir EEX)
  \item Основані на DSL (PureScript, Ocsigen, LiveView, N2O, WebSharper, UrWeb, Flutter)
  \item Клієнтські REST фреймворки з інфраструктурою (vue.js, react.js)
  \item REST мікро-фреймворки
  \item WebSocket фреймворки
\end{itemize}

\newpage

\subsection{Вимоги до безпеки}

Другий підрозділ описує загальні і технічні вимоги до інформаційної безпеки, що діє наскрізно на прикладному,
презентаційному, сесійному і транспортному рівні моделі OSI.

\subsubsection{Загальні вимоги до інформаційної безпеки}

Додати вимоги для персональних даних, приналежність до кластерів.

\begin{enumerate}
  \item Вимоги до інформаційної безпеки повинні відповідати:
    \begin{itemize}
      \item закону України «Про захист інформації в інформаційно-телекомунікаційних системах»;
      \item закону України «Про захист персональних даних»;
      \item закону України «Про основні засади забезпечення кібербезпеки України»;
      \item закону України «Про інформацію»;
      \item правилам забезпечення захисту інформації в інформаційних, телекомунікаційних та інформаційно-телекомунікаційних системах, затвердженим Постановою Кабінету Міністрів України від 29.03.06 № 373;
      \item порядку функціонування електронної системи охорони здоров’я, затвердженому постановою Кабінету Міністрів України від 25.04.2018 № 411;
      \item постанові Кабінету Міністрів України від 21.02.2025 № 205 «Про затвердження Порядку використання засобів інформатизації»;
      \item іншим нормативно-правовим актам, що регулюють питання захисту інформації, у тому числі в інформаційно-телекомунікаційних системах.
    \end{itemize}

  \item Виконавець має забезпечити функціональну можливість засвідчення даних, що вносяться до ERP/1
        користувачами виключно за допомогою кваліфікованого електронного підпису (далі – КЕП):
    \begin{itemize}
      \item в якому наявні атрибути \texttt{certificate-values} та \texttt{revocation-values},
            які свідчать про CAdES-X Long формат довгострокового розширеного підпису КЕП;
      \item що зберігається на захищеному апаратному носії.
    \end{itemize}

  \item Для підключення до ERP/1, виконання функцій формування / перевірки КЕП, для взаємодії з кінцевими
        користувачами в складі додатка мають використовуватись засоби криптографічного захисту інформації,
        які мають чинний позитивний експертний висновок за результатами державної експертизи у сфері
        криптографічного захисту інформації щодо можливості їх використання для криптографічного захисту
        конфіденційної інформації (персональних даних фізичних осіб).

  \item Додаток повинен забезпечувати захист від вразливих компонентів, а саме:
    \begin{itemize}
      \item усі компоненти, включаючи бібліотеки та плагіни, повинні регулярно
            перевірятися на наявність відомих вразливостей, зазначених в базі даних exploit-db (\url{www.exploit-db.com/});
      \item усі застарілі або вразливі компоненти повинні бути оновлені або замінені на безпечні аналоги не пізніше 5 робочих днів з дня виявлення;
      \item система повинна бути налаштована таким чином, щоб уникати використання компонентів, які більше не підтримуються.
    \end{itemize}

  \item Додаток повинен забезпечувати захист від небезпечних налаштувань, а саме:
    \begin{itemize}
      \item всі непотрібні функції та служби повинні бути вимкнені або видалені;
      \item конфігураційні файли повинні бути захищені від несанкціонованого доступу і зберігатися в зашифрованому вигляді.
    \end{itemize}

  \item Додаток повинен забезпечувати захист від порушень управління доступом, а саме:
    \begin{itemize}
      \item доступ до всіх ресурсів повинен здійснюватися згідно з принципом найменших привілеїв;
      \item контроль доступу повинен бути централізованим і надавати права тільки тим користувачам, які мають на це необхідні підстави;
      \item необхідно вести журнал аудиту доступу до конфіденційних даних з можливістю перевірки дій користувачів, термін зберігання яких не менше 3 місяців.
    \end{itemize}

  \item Додаток повинен реалізовувати основні функції захисту:
    \begin{itemize}
      \item забезпечення періодичного резервного копіювання баз даних додатку;
      \item забезпечення однозначної ідентифікації та аутентифікації користувачів додатка;
      \item забезпечення захисту організаційними заходами від несанкціонованого доступу до інформації при її обробці засобами додатка;
      \item контроль за проходженням на мережевому рівні тільки дозволених інформаційних потоків з боку телекомунікаційних мереж до додатку, а також у зворотному напрямку (щонайменше, на рівнях L3, L4 моделі Open Systems Interconnections, OSI);
      \item забезпечення реєстрації подій, пов’язаних з отриманням користувачами доступу до ресурсів додатка (проходження/непроходження автентифікації), зміни налаштувань програмного забезпечення, серверів та комутаційного обладнання;
      \item забезпечення ведення журналів безпеки, усі важливі події безпеки, такі як спроби несанкціонованого доступу, повинні бути зафіксовані у журналах безпеки. Журнали повинні зберігатися не менше 3 місяців, в захищеному вигляді та бути недоступними для модифікації неавторизованими користувачами;
      \item забезпечення конфіденційності та цілісності інформації з обмеженим доступом, що передається каналами зв’язку через незахищене середовище.
    \end{itemize}

  \item Заборонено використовувати проміжні інтерфейси авторизації користувачів, крім веб-сторінки авторизації.

  \item Інформаційна безпека додатка має обов’язково відповідати 10 найбільш важливим перевіркам на безпеку (TOP 10),
        які розміщуються в інформаційному онлайн-документі для розробників і тестувальників з безпеки додатків The
        Open Worldwide Application Security Project (OWASP).
\end{enumerate}

\newpage
\subsubsection{Технічні вимоги до інформаційної безпеки}

\begin{enumerate}
  \item Додаток повинен мати відкритим лише порт 443 для безпечного HTTPS-з'єднання. Інші порти повинні бути закриті.

  \item Додаток повинен використовувати чинний SSL/TLS сертифікат, що видається лише довіреним органом сертифікації.

  \item Додаток повинен використовувати SSL/TLS сертифікат, який відповідає криптографічним стандартам ECC (ECDSA) або RSA (DSA, з 2048-бітними ключами, дозволеними до 2030 року згідно з нормативами NIST).

  \item Додаток повинен забезпечувати захист від SQL-ін'єкцій й використовувати екранування спецсимволів: 
  \texttt{'} (\u0027), \texttt{"} (\u0022), \texttt{\textbackslash} (\u005c), \texttt{--} (\u002d\u002d), \texttt{\%} (\u0025), \texttt{/* */} (\u002f\u002a, \u002a\u002f), \texttt{\#} (\u0023), \texttt{\&} (\u0026), \texttt{\_} (\u005f), \texttt{\textbackslash n} (\u005c\u006e), \texttt{\textbackslash r} (\u005c\u0072), \texttt{\textbackslash t} (\u005c\u0074), \texttt{\%00}.

  \item Додаток повинен забезпечувати захист від вразливостей автентифікації та керування сесіями, а саме:
    \begin{itemize}
      \item сесійні токени повинні зберігатися у cookie з прапорами \texttt{Secure} та \texttt{HttpOnly};
      \item заборона одночасної сесії (сеансу) користувачів на різних пристроях.
    \end{itemize}

  \item Додаток повинен забезпечувати захист від XSS-ін'єкцій, а саме:
    \begin{itemize}
      \item додаток повинен використовувати тільки безпечні заголовки (CSP, HTTPOnly, X-XSS-Protection, X-Frame-Options, HSTS) для запобігання доступу клієнтських сценаріїв до даних;
      \item при внесенні користувачем даних в інтерфейсі додаток повинен попередити користувача та не дозволити
           ввести символи та конструкції (що заповнюються в полях), які можуть бути використані для впровадження
           скриптів або зміни URL, такі як: \texttt{javascript:}, \texttt{vbscript:}, \texttt{data:}, \texttt{onclick},
           \texttt{onload}, \texttt{onsubmit}, \texttt{?} (\u003f), \texttt{'} (\u0027), \texttt{"}
           (\u0022), \texttt{\&} (\u0026), \texttt{<} та \texttt{>} (\u003c \u003e), \texttt{\textbackslash n}
           (\u005c\u006e), \texttt{\textbackslash r} (\u005c\u0072), \texttt{\textbackslash t} (\u005c\u0074),
           \texttt{\%00}, \texttt{/* */} (\u002f\u002a, \u002a\u002f), \texttt{NULL} та \texttt{\#} (\u0023).
    \end{itemize}

  \item Додаток повинен забезпечувати захист від виконання довільного коду (Command Execution), наприклад, \texttt{eval()}, \texttt{exec()}, або \texttt{shell\_exec()}.

  \item Додаток повинен використовувати протокол ``Transport Layer Security (далі – TLS)'' версії не нижче 1.3, що відповідає вимогам чинного законодавства.
\end{enumerate}


\newpage
\subsection{Вимоги до потужності та ємності}

Цей підрозділ визначає процес аналізу та валідації вимог до потужності і ємності нової функціональності,
що є одними з головних технічних ризиків при проєктуванні інформаційних систем.

\begin{enumerate}
  \item Програмний код повинен бути сумісний з горизонтальним масштабуванням процесингових обчислень,
        щоб динамічно збільшувати ємність системи як лінійну залежність від кількості екземплярів процесингових обчислень.
  \item Програмний код має автоматично підтримувати одночасне виконання багатьох невеликих процесів (concurrent)
        та працювати паралельно на кількох ядрах (parallel).
  \item У випадку релізу, який додає функціональність і/або змінює flow по роботі з наявним, слід планувати
        проведення навантажувального тестування з метою визначення впливу на продуктивність системи на середовищі Stage.
  \item Програмне забезпечення, що реалізується, повинно відповідати наступним критеріям:
    \begin{itemize}
      \item Первинне джерело даних, на базі яких проводиться аналіз ризиків потужності та ємності;
      \item Якщо аналітичних даних недостатньо, то прогнозування навантаження повинно відбуватися на основі дотичних сервісів;
      \item Після прийому релізу відділ контролю якості перевіряє виконання вимог до потужності шляхом навантажувального тестування, який передбачає також взаємодію з відділом супроводу і підтримки.
    \end{itemize}

  \item Продуктивність або потужність підсистем визначається наступним набором характеристик (таблиця, яка заповнюється бізнес-аналітиками на основі прогнозованих значень):
    \begin{itemize}
      \item час реакції (активація асинхронної задачі);
      \item пропускна здатність (кількість виконаних задач за одиницю часу);
      \item затримка передачі (latency, час приходу першого байту від асинхронної чи синхронної відповіді).
    \end{itemize}
\end{enumerate}

\subsubsection{Таблиця характеристик продуктивності}

\begin{center}
\begin{tabular}{l l l}
\toprule
\textbf{Характеристика} & \textbf{Опис} & \textbf{Прогноз} \\
\midrule
Час реакції & Активація асинхронної задачі & \textless{} 200 мс \\
Пропускна здатність & Кількість виконаних задач за одиницю часу & 500 задач/с \\
Затримка передачі (latency) & Час приходу першого байту відповіді & \textless{} 150 мс \\
\bottomrule
\end{tabular}
\end{center}

\newpage
\subsection{Вимоги до системи логування}

\subsubsection{Загальні вимоги}

\begin{enumerate}
    \item \textbf{Змістовність.} Кожне повідомлення в логах повинно містити цінну інформацію.
    \item \textbf{Цілісність.} Логи повинні утворювати логічний ланцюжок подій, що призвів до проблеми.
    \item \textbf{Відтворення проблем.} Логи повинні містити достатньо інформації для відтворення проблеми на середовищі розробника.
    \item \textbf{Роздільність.} Кожен сервіс повинен мати власний лог.
    \item \textbf{Лаконічність.} Логи повинні бути читабельними та не містити зайвої інформації.
    \item \textbf{Захист даних.} Приватні та персональні дані (логіни, паролі, секретні ключі, персональні дані клієнтів) не повинні зберігатися в логах в явному вигляді. Використовуйте хешування.
    \item \textbf{Метрики та статистика.} Логи повинні дозволяти побудувати метрики, статистику швидкодії та продуктивності додатка і його компонентів.
    \item \textbf{Історія подій.} Логи повинні містити історію подій до моменту виникнення помилки без потреби запитувати користувачів.
    \item \textbf{Логування помилок.} Логи помилок повинні бути пов’язані з усіма попередніми повідомленнями, що стосуються процесу або запиту. Повідомлення мають бути доступні для пошуку за унікальним ідентифікатором.
    \item \textbf{Логування змін.} Будь-які зміни в конфігурації додатка та його компонентів повинні бути залоговані.
    \item \textbf{Структуровані логи.} Логер повинен бути структурованим і логувати інформацію в форматі JSON.
\end{enumerate}

\subsubsection{Рівні логування}

Рівні логування мають керуватись на рівні конфігурації додатку. Вмикаючи в конфігурації системи один з рівнів логування, також будуть залоговані й повідомлення нижчих рівнів. Наприклад, увімкнувши рівень \texttt{information}, ви матимете записи логів \texttt{warning}, \texttt{error}, \texttt{critical}.

\begin{enumerate}
    \item \textbf{Trace.} Логуються всі властивості об’єктів, параметри методів, виклики методів та виконання кожного рядка коду. Використовується тільки в середовищі розробки/тестування.
    \item \textbf{Debug.} Логуються системні параметри, запити до джерел даних, виклики методів. Використовується на продуктивному середовищі на короткий час для діагностики.
    \item \textbf{Information.} Логуються події, такі як запуск/зупинка компонентів.
    \item \textbf{Warning.} Логуються підозрілі або некритичні позаштатні ситуації.
    \item \textbf{Error.} Логуються події, що переривають поточні операції, але не впливають на подальшу роботу системи.
    \item \textbf{Critical.} Логуються критичні помилки, що перешкоджають подальшій роботі системи (наразі в системі помилка відсутня).
\end{enumerate}

\newpage
\subsubsection{Інформація для вендорів про інструментарії логування}

\begin{enumerate}
    \item \textbf{Filebeat:}
    \begin{itemize}
        \item потрібно використовувати для збору логів з усіх джерел;
        \item повинен бути налаштований на надсилання даних до Elasticsearch.
    \end{itemize}
    \item \textbf{Logstash:}
    \begin{itemize}
        \item має забезпечувати обробку логів з можливістю трансформації, фільтрації та агрегації даних;
        \item конфігурації повинні бути задокументовані та доступні.
    \end{itemize}
    \item \textbf{Elasticsearch:}
    \begin{itemize}
        \item логи повинні зберігатися в Elasticsearch для швидкого пошуку та доступу;
        \item необхідно забезпечити належну масштабованість та продуктивність системи.
    \end{itemize}
    \item \textbf{Kibana:}
    \begin{itemize}
        \item має використовуватися для візуалізації та аналізу збережених логів;
        \item потрібно забезпечити можливість створення графіків та дашбордів для моніторингу.
    \end{itemize}
\end{enumerate}

\subsubsection{Події логування}

\begin{enumerate}
    \item Запуск/зупинка окремих сервісів системи.
    \item Події безпеки типу звернення системи до сервісів системи.
    \item Помилок у роботі системи, таких як комунікаційні, цілісності даних у системі, непередбачувані затримки в обробці інформації.
    \item Критичні події від системи моніторингу (критичний обсяг пам'яті, дискового простору тощо).
    \item Інші події безпеки.
\end{enumerate}

\subsubsection{Структура журналу логування}

Структура журналу логування повинна включати:

\begin{enumerate}
    \item Ідентифікатор запиту в наборах відкритих даних, на який надається відповідь.
    \item Дата та час запиту.
    \item Результат обробки запиту.
    \item Ідентифікатор відповіді від набору відкритих даних, щодо якого надається повідомлення.
    \item Дата та час створення відповіді.
    \item Набір вхідних параметрів.
    \item Час виконання.
\end{enumerate}

\newpage
\subsubsection{Вимоги до структурних логів ELK стеку}

Лог має бути представлений у форматі JSON об’єкта з наступними визначеними полями:

\begin{itemize}
    \item \texttt{time} — log time
    \item \texttt{severity} — log level (\texttt{DEBUG}, \texttt{INFO}, \texttt{WARNING}, \texttt{ERROR})
    \item \texttt{log} — log message
\end{itemize}

Додатково включати по наявності наступні поля, якщо присутні такі метадані:

\begin{itemize}
    \item \texttt{source\_location.file}
    \item \texttt{source\_location.line}
    \item \texttt{source\_location.module}
    \item \texttt{source\_location.function}
    \item \texttt{error.initial\_call} — \texttt{"\#\{module\}.\#\{function\}"} або \texttt{"\#\{module\}.\#\{function\}/\#\{arity\}"}
    \item \texttt{error.reason} — \texttt{"\#\{module\}.\#\{function\}"} або \texttt{"\#\{module\}.\#\{function\}/\#\{arity\}"}
\end{itemize}

Додатково можна включати інші поля за необхідності.

\newpage

\subsection{Вимоги до контролю якості}

\subsubsection{Загальні вимоги до обліку і контролю якості}

Види тестувань:
\begin{itemize}
    \item Внутрішні тести розробників.
    \item Ручне тестування за допомогою Postman, curl, wscat, тощо.
    \item Автоматичне регресивне тестування з використанням кваліфікаційного електронного підпису.
    \item Навантажувальне тестування.
\end{itemize}

Тестувальна документація має трактувати кроки, терміни та результати тестування.
Склад документації з тестування залежить від конкретної задачі, але повинен містити:

\begin{itemize}
    \item \textbf{Об'єкт тестування.} Необхідно чітко визначити, який об'єкт тестування підлягає приймальному тестуванню (модуль, програмне забезпечення, веб-додаток, мобільний додаток або комплексна система). Вказати конкретну функціональність та компоненти.
    \item \textbf{Критерії успішності.} Критерії, які повинні бути виконані для визнання тестування успішним (функціональні вимоги, продуктивність, стабільність, безпека, сумісність тощо).
    \item \textbf{Сценарій тестування.} Опис тестового середовища, набір сценаріїв тестування та покрокових процедур (тест-кейси).
    \item \textbf{Умови тестування.} Конфігурації системи, вхідні дані, налаштування, середовища виконання тощо.
    \item \textbf{План тестування.} Документ, що описує весь обсяг робіт із тестування.
    \item \textbf{Критерії прийняття.} Кількість та тип прийнятних дефектів, виконання функціональностей, бізнес-сценаріїв тощо.
    \item \textbf{Протоколи результатів тестування.} Загальна оцінка (успішно/неуспішно), вплив тестового середовища, результати кожного тест-кейсу.
    \item \textbf{Стратегія контролю версій.} Контроль та управління версіями документів, які використовуються в тестуванні.
    \item \textbf{Вимоги до навантаження та продуктивності.} Обсяги даних, кількість одночасних користувачів, транзакцій, очікувані показники продуктивності та час відповіді.
    \item \textbf{Звітність та документування.} Формат звітів, структура, висновки та рекомендації.
\end{itemize}

\newpage
\subsubsection{Вимоги до ручного тестування}

Ручне тестування повинно оцінювати функціональність з точки зору звичайного
користувача та надавати максимально розгорнутий і зрозумілий звіт.
Необхідно виконати перевірку:

\begin{itemize}
    \item форм користувача (коректне заповнення, виведення помилок, обов’язкові/необов’язкові поля, календарі, меню тощо);
    \item посилань та навігації сайту;
    \item форм реєстрації та авторизації;
    \item роботи з базою даних (додавання, редагування, видалення даних, завантаження файлів);
    \item файлів cookie;
    \item пошукового рядка.
\end{itemize}

Необхідно виконати перевірку \textbf{нетипових сценаріїв} для виявлення несподіваних помилок і багів.

\textbf{Перевірка безпеки:}
\begin{itemize}
    \item Конфіденційність (обмеження доступу до особистих даних);
    \item Цілісність (можливість відновлення інформації після атаки);
    \item Доступність (розмежування рівнів доступу).
\end{itemize}

\textbf{Оцінка дизайну:}
\begin{itemize}
    \item Простота експлуатації;
    \item Зручність навігації;
    \item Відповідність контенту (шрифти, кольори, розташування елементів);
    \item Мінімальна кількість кліків для виконання операцій.
\end{itemize}

\textbf{Тестування на сумісність:}
\begin{itemize}
    \item На різних пристроях (ПК, планшет, телефон) та ОС (Windows, macOS, Android);
    \item У різних браузерах (Edge, Firefox, Chrome, Safari, Opera).
\end{itemize}

\textbf{Тестування продуктивності}
\begin{itemize}
    \item Навантажувальне тестування;
    \item Стрес-тестування;
    \item Об’ємне тестування;
    \item Тестування надійності.
\end{itemize}

Необхідно виконати \textbf{регресивне тестування} після внесення змін у код.
Функціональні вимоги до HTTP API тестування:

\begin{itemize}
    \item Правильність відповідей та використання коректних HTTP-статус-кодів
    \item Валідація вхідних даних (формати, обов’язкові поля, медичні коди)
    \item Захист від помилок та чіткі повідомлення про помилки
    \item Правильний формат даних у відповідях
    \item Перевірка обов’язкових полів.
\end{itemize}

Вимоги до специфікацій технічної документації:
Технічна документація повинна містити:

\begin{itemize}
    \item Архітектурний опис системи та взаємодії компонентів
    \item Опис модулів та функцій
    \item API-специфікацію (ендпоінти, методи, запити, відповіді, статуси)
    \item Опис бази даних (схема, таблиці, зв’язки);
    \item Опис інтерфейсу користувача (wireframes, прототипи, адаптивність)
\end{itemize}

\subsubsection{Вимоги до формату сценаріїв та протоколів тестування}

Сценарії надаються у вигляді таблиці У полях «Від замовника» та «Від виконавця» вказується ПІБ.
Зазначаються дати надання та виконання сценаріїв.
Протокол тестування формується у вигляді звіту Jira або у форматі XLS/DOCX.
Протокол повинен містити:

\begin{itemize}
    \item Найменування системи та номер версії;
    \item Назву тестового середовища;
    \item Перелік виконаних тест-кейсів.
\end{itemize}

Для кожного тест-кейсу вказується:

\begin{itemize}
    \item Назва тест-кейсу;
    \item Статус;
    \item Дата виконання;
    \item ПІБ тестувальника;
    \item Перелік незакритих дефектів;
    \item Підсумковий висновок.
\end{itemize}

\begin{table}[htbp]
\centering
\caption{Шаблон протоколу тестування}
\label{tab:protocol_template}
\begin{tabular}{|l|l|}
\hline
\textbf{ID сценарію} & \\
\hline
\textbf{Назва сценарію} & \\
\hline
\textbf{ID тестового випадку} & \\
\hline
\textbf{Передумови} & \\
\hline
\textbf{Кроки відтворення / Умови} & \\
\hline
\textbf{Очікуваний результат} & \\
\hline
\textbf{Статус} & \\
\hline
\textbf{Коментар} & \\
\hline
\textbf{Дата тестування} & \\
\hline
\textbf{Тестувальник} & \\
\hline
\textbf{Від замовника} & \\
\hline
\textbf{Від виконавця} & \\
\hline
\end{tabular}
\end{table}

\newpage


\subsubsection{Вимоги до плану тестування}

Перед початком тестування функціональності має бути підготовлено план тестування.
План тестування — це документ, в якому визначені об’єми, ресурси та календарний
план тестування. Визначається відділом контролю якості.
План тестування повинен містити наступну інформацію:

\begin{itemize}
    \item \textbf{Об'єкт тестування.} Назва проекту або компонента, який підлягає тестуванню.
    \item \textbf{Мету тестування.} Виявлення дефектів, перевірка функціональності, валідація вимог, оцінка продуктивності тощо.
    \item \textbf{Стратегія тестування.} Типи тестів: модульне, інтеграційне, системне, регресійне, продуктивності, безпеки тощо.
    \item \textbf{Обсяг тестування.} Кількість тестових сценаріїв, наборів даних, ітерацій, тривалість тестового циклу.
    \item \textbf{Ресурси тестування.} Людські ресурси, апаратне та програмне забезпечення, тестові середовища, тестові дані.
    \item \textbf{Графік тестування.} Дати початку та закінчення тестового циклу, терміни проведення окремих видів тестів.
    \item \textbf{Ризики тестування.} Ідентифікація, оцінка ризиків та заходи щодо їх зниження.
    \item \textbf{Критерії прийняття.} Умови прийняття системи після тестування.
    \item \textbf{Відповідальності та ролі.} Ролі та відповідальності учасників тестування (тестувальники, керівники проекту, розробники тощо).
\end{itemize}

\subsubsection{Вимоги до автоматичного тестування}

Вся функціональність сервісу повинна бути повністю покрита автоматичними тестами.

Виконавець зобов’язаний надати документацію, яка містить:
\begin{itemize}
    \item Об’єкт автотестування;
    \item Мету автотестування;
    \item Типи тестів та сценарії;
    \item Використовувані фреймворки та інструменти;
    \item Критерії успішності;
    \item Інтеграцію в процес CI/CD;
    \item Описи тестів, підготовку тестових даних, передумови та очікувані результати.
\end{itemize}




\subsubsection{Вимоги до результатів приймального (UAT) тестування}

Приймальне (UAT) тестування має бути виконано на Demo-середовищі (якщо інше не погоджено з Замовником).
Замовник затверджує висновок про відповідність функціональності вимогам технічного завдання згідно з таблицею 5.5.1.
Після успішного UAT Виконавець передає поставку програмного забезпечення для інсталяції та налаштування на середовищах Замовника.
Прийняття функціональності можливе повністю або частинами згідно з календарним планом.
Функціональність має бути прийнята Замовником в повному обсязі згідно з календарним планом виконання робіт.

\newpage
\subsection{Вимоги до інтерфейсів}

Даний підрозділ визначає ергономічні та технічні вимоги до інтерфейсів користувача адміністративних панелей.

\subsubsection{Загальні вимоги}

Інтерфейс користувача адміністративної панелі є критичним елементом для ефективного управління та
адміністрування програмного забезпечення чи веб-сайту.

\begin{itemize}
    \item \textbf{Інтуїтивність та простота.} Інтерфейс повинен бути легким для використання, навіть для користувачів без глибоких технічних знань.
    \item \textbf{Логічна структура та навігація.} Забезпечення логічної організації елементів та легкої навігації між різними функціями та розділами.
    \item \textbf{Мобільна сумісність.} Адаптивний дизайн для коректного відображення на різних пристроях, включаючи мобільні телефони та планшети.
    \item \textbf{Гнучкість налаштувань.} Можливість налаштування інтерфейсу відповідно до унікальних потреб та завдань адміністратора.
    \item \textbf{Зручне управління даними.} Легкість управління, внесення та вилучення даних з адміністративної панелі.
    \item \textbf{Інформативність.} Забезпечення користувача інформацією про стан системи, статистику та повідомлення про можливі проблеми.
    \item \textbf{Безпека та авторизація.} Захист інтерфейсу від несанкціонованого доступу, використання протоколів шифрування та можливість налаштування рівнів доступу.
    \item \textbf{Інтеграція з іншими системами.} Можливість інтеграції з іншими системами та сервісами.
    \item \textbf{Спрощений процес розгортання та оновлення.} Легкість встановлення та оновлення адмінпанелі для мінімізації перерв у роботі.
    \item \textbf{Підтримка різних мов.} Locaize (продукт який використовується в кабінеті пацієнта).
    \item \textbf{Підтримка технічних засобів.} Забезпечення сумісності з різними веб-браузерами та операційними системами.
\end{itemize}

\subsubsection{Вимоги до відображення інформації}

\begin{itemize}
    \item \textbf{Кросбраузерність.} Система має бути адаптованою для використання в усіх сучасних браузерах.
    \item \textbf{Адаптація для людей з вадами зору.} Інтерфейс має бути адаптований для людей з порушенням зору згідно з ДСТУ EN 301 549:2022 та рекомендацій WCAG 2.1 (W3C Web Content Accessibility Guidelines):
    \begin{itemize}
        \item контрастні кольори;
        \item вибір розміру шрифту.
    \end{itemize}
    \item \textbf{Адаптація до розмірів екрана на пристроях.} Веб-сайт має бути адаптивним і коректно масштабуватися на екранах різного розміру.
\end{itemize}

\newpage
\subsubsection{Вимоги до швидкодії додатка}

Швидкість завантаження та відображення табличних частин.
Час повного завантаження таблиці з 2000 рядками повинен становити не більше 1 с,
а відображення видимої частини повинна становити не більше 300 мс.

\subsubsection{Вимоги до технології реалізації додатку}

\begin{itemize}
    \item \textbf{Single Page Application.} Використання наперед сформованих HTML-частин на стороні сервера і модифікація структури DOM-елементів в браузері відносно WebSocket.
    \item \textbf{Data-on-Wire.} Використання тільки даних між клієнтом і сервером, що вимагають розвиненої інфраструктури (шаблонізаторів і клієнтської шини повідомлень для графічних контрольних елементів) і мов програмування для об'ємних клієнтів.
    \item \textbf{Figma.} Артефакти додатку (assets) та/або демо-прототип UI/UX повинні надаватися в Figma.
\end{itemize}

\newpage
\subsection{Юридичні вимоги}

\subsubsection{Зміст}

  \begin{itemize}
\item Перелік юридичних документів, що входять до комплекту поставки та надаються в паперовому вигляді
  \begin{enumerate}
\item Договір про передання (відчуження) майнових прав інтелектуальної власності
\item Акт прийняття-передання наданих послуг
\item Акт приймання-передачі майнових прав на Комп'ютерну програму
\item Договори з субпідрядниками (за наявності)
\item Акт прийняття-передання по електронному носії інформації (опис електронного носія інформації)
  \end{enumerate}
\item Вимоги до патентної чистоти
\item Вимоги до розвитку та модернізації функціоналу
\item Гарантійна (технічна) підтримка функціоналу
\item Перелік юридичних документів
\item Вимоги до передачі прав на інтелектуальну власність
  \end{itemize}

На момент момент укладення договору щодо розроблення Системи Виконавець має врегулювати усі відносини з авторами та третіми особами, які залучаються ним до створення Системи, а також надати підтвердження виникнення у нього в майбутньому майнових прав на Систему та її компоненти:

1. Надати підтвердження укладення з авторами Системи та її компонентів договорів, предметом яких є створення на замовлення об’єктів права інтелектуальної власності. Умовами договорів має бути врегульовано питання щодо: (1) отримання замовлення (завдання) щодо створення Системи та її компонентів; (2) переходу права власності на результати виконаних робіт; (3) отримання авторської винагороди (у т. ч. паушальної) у зв’язку з передачею прав щодо використання без часових, предметних, кількісних та територіальних обмежень; (4) компенсацій та інших платежів за передачу майнових прав; (5) немайнових прав автора, зокрема права на внесення змін до твору; (6) оригінальності твору або отримання відповідних дозволів щодо створення похідних творів. Якщо Система та її компоненти будуть створюватися як службовий твір - надати підтвердження трудових відносин з авторами. 

2. Врегулювати питання щодо оригінальності твору, а саме створення Системи та її компонентів з використанням інших об’єктів права інтелектуальної власності, які раніше не створювалися, або щодо отримання всіх прав на таке використання.

3. Врегулювати питання щодо відсутності будь-яких майнових чи немайнових вимог з боку інших осіб, для яких зазначена Системи та її компоненти, створювалася раніше, або яким передавалися права на таку Систему та її компонентів, та забезпечити умови ексклюзивності.

4. Врегулювати інші питання, важливі для забезпечення переходу до держави в особі ДСА України всіх прав на Систему та її компоненти  без часових, предметних, кількісних, територіальних чи інших обмежень.

\newpage
\section*{Висновок}

Дотримання принципу «політики → вимоги → технічне завдання» та обов’язкове використання ДСТУ гарантує технологічну незалежність, аудитованість та відповідність державним нормам. Конкретні продукти обираються лише на етапі технічного завдання.

\begin{thebibliography}{9}

\bibitem{ecoz:rehab} Міністерство охорони здоров’я України. \newblock Технічні вимоги Реабілітація 2.0. \newblock Версія 08.08.2024.
\bibitem{dstu:3973} ДСТУ 3973-2000. \newblock Система розроблення та постачання продукції на виробництво. Правила виконання науково-дослідних робіт. Загальні положення.
\bibitem{dstu:3974} ДСТУ 3974-2000. \newblock Система розроблення та постачання продукції на виробництво. Правила виконання науково-конструкторських робіт. Загальні положення.
\bibitem{dstu:42010} ДСТУ 42010:2018. \newblock Інженерія систем і програмних засобів. Опис архітектури.
\bibitem{law:info-protection} Закон України «Про захист інформації в інформаційно-телекомунікаційних системах».
\bibitem{law:personal-data} Закон України «Про захист персональних даних».
\bibitem{law:cybersecurity} Закон України «Про основні засади забезпечення кібербезпеки України».
\bibitem{owasp:top10} OWASP Foundation. \newblock OWASP Top 10: The Ten Most Critical Web Application Security Risks. \newblock \url{https://owasp.org/www-project-top-ten/}.
\bibitem{nist:sp800-57} NIST Special Publication 800-57 Part 1 Revision 5. \newblock Recommendation for Key Management.
\bibitem{iso:42010} ISO/IEC/IEEE 42010:2022. \newblock Systems and software engineering — Architecture description.
\bibitem{zachman} Zachman, J.A. \newblock A Framework for Information Systems Architecture. \newblock IBM Systems Journal, 1987.

\end{thebibliography}

\end{document}
